% =========================================================
% Chapter: Temperature, Heat, and Thermal Expansion
% POSN 1 Physics Preparation Notes
% Language: Thai
% Compiler: XeLaTeX
% =========================================================

\chapter{อุณหภูมิ ความร้อน และการขยายตัวเชิงความร้อน}

\begin{center}
    {\Large \textbf{Temperature, Heat, and Thermal Expansion}}\\[4pt]
    {\normalsize เอกสารเตรียมสอบคัดเลือก สอวน. ฟิสิกส์ ค่าย 1}\\[6pt]
    {\normalsize จัดทำโดย \textbf{Tames Thanakrit Damduan}}\\[2pt]
    {\footnotesize \textcopyright\ 2026 Tames Thanakrit Damduan. All rights reserved.}\\[8pt]
\end{center}

\section*{เป้าหมายของบทเรียน}

หลังจากเรียนบทนี้ นักเรียนควรทำสิ่งต่อไปนี้ได้

\begin{enumerate}
    \item เข้าใจความหมายของอุณหภูมิ สมดุลความร้อน และกฎข้อที่ศูนย์ของเทอร์โมไดนามิกส์ในระดับที่จำเป็น
    \item แปลงหน่วยอุณหภูมิระหว่างองศาเซลเซียสและเคลวินได้
    \item ใช้สมการการขยายตัวเชิงเส้น เชิงพื้นที่ และเชิงปริมาตรได้
    \item วิเคราะห์รูหรือช่องว่างในวัตถุที่ขยายตัวเมื่ออุณหภูมิเปลี่ยนได้
    \item ใช้ความจุความร้อนจำเพาะและหลักอนุรักษ์พลังงานในโจทย์แคลอริเมตรีได้
    \item แยกช่วงการเปลี่ยนอุณหภูมิออกจากช่วงเปลี่ยนสถานะได้
    \item ใช้สมการการนำความร้อน ฟลักซ์ความร้อน และความต้านทานความร้อนได้
    \item วิเคราะห์ท่อนำความร้อนหลายท่อนต่ออนุกรมและหาอุณหภูมิที่รอยต่อได้
    \item เชื่อมโยงการขยายตัวเชิงความร้อนกับโจทย์ลูกตุ้มและสถานการณ์ระดับ สอวน. ได้
\end{enumerate}

\begin{tcolorbox}[title=แนวคิดสำคัญของบทนี้]
บทนี้มีแกนหลักสองส่วน

\[
\text{พลังงานความร้อนที่รับหรือคาย}
\]

และ

\[
\text{ผลของอุณหภูมิต่อขนาดและการไหลของพลังงาน}
\]

โจทย์ที่ยากมักไม่ได้ใช้สูตรเดียว แต่ต้องแบ่งสถานการณ์เป็นช่วง แล้วเลือกกฎอนุรักษ์พลังงานหรือเงื่อนไขอัตราการไหลคงตัวให้ถูกต้อง
\end{tcolorbox}

\newpage{}

% =========================================================
\section{อุณหภูมิและสมดุลความร้อน}

\subsection{อุณหภูมิคืออะไร}

อุณหภูมิเป็นปริมาณที่ใช้บอกสภาวะเชิงความร้อนของระบบ ถ้าวัตถุสองชิ้นมีอุณหภูมิต่างกันและสัมผัสกัน พลังงานความร้อนจะถ่ายโอนจากวัตถุที่มีอุณหภูมิสูงกว่าไปสู่วัตถุที่มีอุณหภูมิต่ำกว่า จนกระทั่งทั้งสองอยู่ในสมดุลความร้อน

\begin{tcolorbox}[title=สมดุลความร้อน]
วัตถุสองชิ้นอยู่ในสมดุลความร้อนเมื่อไม่มีการถ่ายโอนพลังงานความร้อนสุทธิระหว่างกัน

ในสภาวะนี้

\[
T_1=T_2
\]
\end{tcolorbox}

\subsection{กฎข้อที่ศูนย์ของเทอร์โมไดนามิกส์}

ถ้าระบบ $A$ อยู่ในสมดุลความร้อนกับระบบ $C$ และระบบ $B$ อยู่ในสมดุลความร้อนกับระบบ $C$ แล้วระบบ $A$ และ $B$ จะอยู่ในสมดุลความร้อนซึ่งกันและกัน

\[
T_A=T_C
\]

และ

\[
T_B=T_C
\]

ดังนั้น

\[
T_A=T_B
\]

หลักนี้ทำให้เราสามารถใช้เทอร์โมมิเตอร์เป็นระบบตัวกลางในการวัดอุณหภูมิได้

\subsection{องศาเซลเซียสและเคลวิน}

ความแตกต่างของอุณหภูมิหนึ่งองศาเซลเซียสมีขนาดเท่ากับความแตกต่างหนึ่งเคลวิน

\[
\Delta T_{\si{\degreeCelsius}}=\Delta T_{\si{K}}
\]

แต่ค่าศูนย์ของสเกลต่างกัน

\[
T_{\si{K}}=T_{\si{\degreeCelsius}}+273.15
\]

ในโจทย์คำนวณระดับค่าย 1 มักใช้

\[
T_{\si{K}}\approx T_{\si{\degreeCelsius}}+273
\]

\begin{tcolorbox}[title=ข้อควรระวัง]
ในสมการที่ใช้เฉพาะผลต่างอุณหภูมิ เช่น

\[
Q=mc\Delta T
\]

สามารถใช้ผลต่างในหน่วยองศาเซลเซียสหรือเคลวินได้

แต่ในสมการที่ใช้อุณหภูมิสัมบูรณ์ เช่น การแผ่รังสีความร้อน ต้องใช้อุณหภูมิหน่วยเคลวิน
\end{tcolorbox}

% =========================================================
\section{การขยายตัวเชิงความร้อน}

\subsection{การขยายตัวเชิงเส้น}

เมื่ออุณหภูมิของแท่งวัตถุเปลี่ยนไปเล็กน้อย ความยาวของแท่งจะเปลี่ยนโดยประมาณเป็นสัดส่วนกับความยาวเดิมและผลต่างอุณหภูมิ

\[
\Delta L=\alpha L_0\Delta T
\]

ดังนั้นความยาวใหม่คือ

\[
L=L_0+\Delta L
\]

\[
L=L_0(1+\alpha\Delta T)
\]

โดยที่

\[
\alpha=\text{สัมประสิทธิ์การขยายตัวเชิงเส้น}
\]

หน่วยของ $\alpha$ คือ

\[
\si{K^{-1}}
\]

หรือ

\[
\si{\degreeCelsius^{-1}}
\]

\begin{tcolorbox}[title=การขยายตัวเชิงเส้น]
\[
\boxed{\Delta L=\alpha L_0\Delta T}
\]

\[
\boxed{L=L_0(1+\alpha\Delta T)}
\]
\end{tcolorbox}

\begin{examplebox}{ตัวอย่าง: ช่องว่างระหว่างราง}
รางโลหะแต่ละท่อนยาว $12.0\,\si{m}$ ที่อุณหภูมิ $20\,\si{\degreeCelsius}$ โลหะมี

\[
\alpha=1.2\times 10^{-5}\,\si{K^{-1}}
\]

ถ้าอุณหภูมิสูงสุดเป็น $60\,\si{\degreeCelsius}$ จงหาช่องว่างต่ำสุดที่ต้องเว้นระหว่างรางสองท่อน เพื่อไม่ให้รางดันกัน

ผลต่างอุณหภูมิคือ

\[
\Delta T=60-20=40\,\si{K}
\]

การเพิ่มความยาวของรางหนึ่งท่อนคือ

\[
\Delta L=\alpha L_0\Delta T
\]

\[
\Delta L=(1.2\times 10^{-5})(12.0)(40)
\]

\[
\Delta L=5.76\times 10^{-3}\,\si{m}
\]

ดังนั้นควรเว้นช่องว่างอย่างน้อย

\[
\boxed{5.76\,\si{mm}}
\]
\end{examplebox}

\subsection{รูในแผ่นโลหะขยายหรือหด}

ถ้าแผ่นโลหะมีรูและถูกให้ความร้อน รูจะขยายตัวเหมือนกับว่ารูนั้นทำจากวัสดุชนิดเดียวกับแผ่นโลหะ

\begin{tcolorbox}[title=หลักของรูในวัตถุ]
เมื่อวัตถุขยายตัวอย่างสม่ำเสมอ ทุกระยะเชิงเส้นจะถูกคูณด้วยตัวประกอบเดียวกัน

\[
1+\alpha\Delta T
\]

ดังนั้นเส้นผ่านศูนย์กลางของรูจะเพิ่มขึ้นเมื่อให้ความร้อน
\end{tcolorbox}

\subsection{การขยายตัวเชิงพื้นที่}

ให้แผ่นสี่เหลี่ยมมีด้านเดิม

\[
L_0
\]

และ

\[
W_0
\]

พื้นที่เดิมคือ

\[
A_0=L_0W_0
\]

เมื่ออุณหภูมิเพิ่มขึ้น $\Delta T$

\[
L=L_0(1+\alpha\Delta T)
\]

\[
W=W_0(1+\alpha\Delta T)
\]

ดังนั้น

\[
A=A_0(1+\alpha\Delta T)^2
\]

กระจายพจน์

\[
A=A_0\left(1+2\alpha\Delta T+\alpha^2\Delta T^2\right)
\]

สำหรับการขยายตัวเล็กน้อย

\[
\alpha\Delta T\ll 1
\]

จึงละพจน์กำลังสองได้

\[
A\approx A_0(1+2\alpha\Delta T)
\]

ดังนั้น

\[
\Delta A\approx 2\alpha A_0\Delta T
\]

\begin{tcolorbox}[title=การขยายตัวเชิงพื้นที่]
สำหรับของแข็งที่ขยายตัวเท่ากันทุกทิศทาง

\[
\boxed{\Delta A\approx 2\alpha A_0\Delta T}
\]
\end{tcolorbox}

\subsection{การขยายตัวเชิงปริมาตร}

ในทำนองเดียวกัน ถ้าของแข็งมีมิติเดิมสามแกน และแต่ละแกนเพิ่มด้วยตัวประกอบ

\[
1+\alpha\Delta T
\]

จะได้

\[
V=V_0(1+\alpha\Delta T)^3
\]

กระจายพจน์

\[
V=V_0\left(1+3\alpha\Delta T+3\alpha^2\Delta T^2+\alpha^3\Delta T^3\right)
\]

เมื่อ

\[
\alpha\Delta T\ll 1
\]

ละพจน์อันดับสูงได้

\[
V\approx V_0(1+3\alpha\Delta T)
\]

ดังนั้น

\[
\Delta V\approx 3\alpha V_0\Delta T
\]

ถ้าใช้สัมประสิทธิ์การขยายตัวเชิงปริมาตร $\beta$

\[
\Delta V=\beta V_0\Delta T
\]

จึงได้โดยประมาณว่า

\[
\beta\approx 3\alpha
\]

\begin{tcolorbox}[title=การขยายตัวเชิงปริมาตร]
\[
\boxed{\Delta V=\beta V_0\Delta T}
\]

สำหรับของแข็งที่ขยายตัวเท่ากันทุกทิศทาง

\[
\boxed{\beta\approx 3\alpha}
\]
\end{tcolorbox}

\subsection{หลักการประมาณสำหรับปริมาณที่เปลี่ยนเล็กน้อย}

ถ้า

\[
|\epsilon|\ll 1
\]

เรามักใช้

\[
(1+\epsilon)^n\approx 1+n\epsilon
\]

ตัวอย่างเช่น

\[
\sqrt{1+\epsilon}
=
(1+\epsilon)^{1/2}
\approx
1+\frac{1}{2}\epsilon
\]

หลักนี้สำคัญในโจทย์การขยายตัวเชิงความร้อนและโจทย์ลูกตุ้ม

\begin{exercisebox}{แบบฝึกหัด 1: การขยายตัวเชิงความร้อน}
\begin{enumerate}
    \item แท่งโลหะยาว $2.0\,\si{m}$ มี $\alpha=1.5\times10^{-5}\,\si{K^{-1}}$ เมื่ออุณหภูมิเพิ่มขึ้น $80\,\si{K}$ ความยาวเพิ่มขึ้นเท่าไร
    \item แผ่นโลหะมีรูวงกลมเส้นผ่านศูนย์กลาง $4.00\,\si{cm}$ ที่อุณหภูมิเดิม ถ้า $\alpha\Delta T=2.0\times10^{-3}$ เส้นผ่านศูนย์กลางใหม่ของรูเป็นเท่าไร
    \item ลูกบาศก์โลหะมีปริมาตร $V_0$ และสัมประสิทธิ์การขยายตัวเชิงเส้น $\alpha$ จงหาปริมาตรใหม่โดยละพจน์อันดับสองขึ้นไป
    \item เหตุใดการเทน้ำร้อนบนฝาขวดโลหะจึงช่วยเปิดฝาได้ง่ายขึ้น
\end{enumerate}
\end{exercisebox}

% =========================================================
\section{ความร้อน ความจุความร้อน และความจุความร้อนจำเพาะ}

\subsection{ความร้อนเป็นพลังงานที่ถ่ายโอน}

ความร้อน $Q$ ไม่ใช่สิ่งที่วัตถุ ``มีอยู่'' แบบเดียวกับอุณหภูมิ แต่เป็นพลังงานที่ถ่ายโอนเนื่องจากความแตกต่างของอุณหภูมิ

\[
Q>0
\]

เมื่อระบบรับพลังงานความร้อน และ

\[
Q<0
\]

เมื่อระบบคายพลังงานความร้อน

\subsection{ความจุความร้อน}

ความจุความร้อน $C$ ของวัตถุคือพลังงานความร้อนที่ต้องใช้เพื่อให้อุณหภูมิของวัตถุเพิ่มขึ้นหนึ่งหน่วย

\[
C=\frac{Q}{\Delta T}
\]

ดังนั้น

\[
Q=C\Delta T
\]

\subsection{ความจุความร้อนจำเพาะ}

ความจุความร้อนจำเพาะ $c$ เป็นความจุความร้อนต่อหนึ่งหน่วยมวล

\[
c=\frac{C}{m}
\]

ดังนั้น

\[
C=mc
\]

และ

\[
Q=mc\Delta T
\]

\begin{tcolorbox}[title=ความร้อนที่ทำให้อุณหภูมิเปลี่ยน]
\[
\boxed{Q=mc\Delta T}
\]

สมการนี้ใช้เมื่อสารยังอยู่ในสถานะเดิมตลอดช่วงที่พิจารณา
\end{tcolorbox}

\subsection{กราฟอุณหภูมิกับพลังงานความร้อน}

ถ้าสารยังไม่เปลี่ยนสถานะ

\[
Q=mc\Delta T
\]

จัดรูปได้

\[
\Delta T=\frac{Q}{mc}
\]

ดังนั้นความชันของกราฟ $T$ กับ $Q$ คือ

\[
\frac{\Delta T}{Q}=\frac{1}{mc}
\]

วัตถุที่มี $mc$ มาก จะเปลี่ยนอุณหภูมิช้ากว่าเมื่อรับพลังงานเท่ากัน

\begin{examplebox}{ตัวอย่าง: ผสมของเหลวชนิดเดียวกัน}
นำน้ำมวล $m$ ที่อุณหภูมิ $80\,\si{\degreeCelsius}$ ผสมกับน้ำมวล $2m$ ที่อุณหภูมิ $20\,\si{\degreeCelsius}$ ในภาชนะฉนวน จงหาอุณหภูมิสุดท้าย

ความร้อนที่น้ำร้อนคายออกคือ

\[
Q_{\text{hot}}=mc(80-T_f)
\]

ความร้อนที่น้ำเย็นรับคือ

\[
Q_{\text{cold}}=2mc(T_f-20)
\]

เพราะภาชนะเป็นฉนวน

\[
Q_{\text{hot}}=Q_{\text{cold}}
\]

ดังนั้น

\[
mc(80-T_f)=2mc(T_f-20)
\]

ตัด $mc$

\[
80-T_f=2T_f-40
\]

\[
120=3T_f
\]

\[
\boxed{T_f=40\,\si{\degreeCelsius}}
\]
\end{examplebox}

% =========================================================
\section{แคลอริเมตรีและหลักอนุรักษ์พลังงาน}

\subsection{ระบบปิดเชิงความร้อน}

ถ้าภาชนะเป็นฉนวน และไม่มีงานจากภายนอกที่ต้องนำมาคิด พลังงานความร้อนที่วัตถุหนึ่งคายออกจะเท่ากับพลังงานความร้อนที่วัตถุอื่นรับเข้า

\[
\sum Q=0
\]

หรือเขียนเป็นภาษาง่าย ๆ ว่า

\[
\text{ความร้อนที่คายออก}
=
\text{ความร้อนที่รับเข้า}
\]

\begin{tcolorbox}[title=หลักของโจทย์ผสมสาร]
\[
\boxed{\sum_i Q_i=0}
\]

กำหนดเครื่องหมายอย่างสม่ำเสมอ หรือแยกเป็น

\[
\boxed{Q_{\text{lost}}=Q_{\text{gained}}}
\]
\end{tcolorbox}

\subsection{ขั้นตอนแก้โจทย์แคลอริเมตรี}

\begin{enumerate}
    \item ระบุวัตถุทุกชิ้นที่รับหรือคายพลังงาน
    \item ตรวจสอบว่ามีการเปลี่ยนสถานะหรือไม่
    \item แบ่งการเปลี่ยนแปลงเป็นช่วง
    \item เขียน $Q$ ของแต่ละช่วง
    \item ใช้ $\sum Q=0$
    \item ตรวจสอบว่าอุณหภูมิสุดท้ายที่ได้สอดคล้องกับสมมติฐานหรือไม่
\end{enumerate}

\begin{examplebox}{ตัวอย่าง: โลหะร้อนใส่ลงในน้ำ}
โลหะมวล $0.20\,\si{kg}$ มีความจุความร้อนจำเพาะ

\[
c_m=450\,\si{J/(kg.K)}
\]

เริ่มที่ $180\,\si{\degreeCelsius}$ นำไปใส่ในน้ำมวล $0.30\,\si{kg}$ ที่ $20\,\si{\degreeCelsius}$ โดยให้

\[
c_w=4200\,\si{J/(kg.K)}
\]

และภาชนะเป็นฉนวน จงหาอุณหภูมิสุดท้าย

โลหะคายความร้อน

\[
Q_m=(0.20)(450)(180-T_f)
\]

น้ำรับความร้อน

\[
Q_w=(0.30)(4200)(T_f-20)
\]

ตั้งสมการ

\[
(0.20)(450)(180-T_f)
=
(0.30)(4200)(T_f-20)
\]

\[
90(180-T_f)=1260(T_f-20)
\]

\[
16200-90T_f=1260T_f-25200
\]

\[
41400=1350T_f
\]

ดังนั้น

\[
\boxed{T_f\approx 30.7\,\si{\degreeCelsius}}
\]
\end{examplebox}

% =========================================================
\section{การเปลี่ยนสถานะและความร้อนแฝง}

\subsection{ความร้อนแฝง}

ระหว่างการเปลี่ยนสถานะ อุณหภูมิของสารอาจคงที่แม้ยังมีการถ่ายโอนพลังงาน ความร้อนส่วนนี้ใช้ในการเปลี่ยนโครงสร้างระดับจุลภาคของสาร

\[
Q=mL
\]

โดยที่

\[
L=\text{ความร้อนแฝงจำเพาะ}
\]

สำหรับการหลอมเหลว ใช้

\[
L_f
\]

สำหรับการกลายเป็นไอ ใช้

\[
L_v
\]

\begin{tcolorbox}[title=ความร้อนแฝง]
\[
\boxed{Q=mL}
\]

ระหว่างการเปลี่ยนสถานะที่ความดันคงที่ อุณหภูมิจะคงที่จนกว่าการเปลี่ยนสถานะจะเสร็จ
\end{tcolorbox}

\subsection{การแบ่งช่วงของโจทย์น้ำแข็ง}

ถ้าน้ำแข็งเริ่มที่อุณหภูมิต่ำกว่า $0\,\si{\degreeCelsius}$ และสุดท้ายกลายเป็นน้ำที่อุณหภูมิสูงกว่า $0\,\si{\degreeCelsius}$ ต้องแบ่งเป็นสามช่วง

\[
\text{ช่วงที่ 1: ทำให้น้ำแข็งอุ่นขึ้นถึง }0\,\si{\degreeCelsius}
\]

\[
Q_1=m c_{\text{ice}}(0-T_i)
\]

\[
\text{ช่วงที่ 2: หลอมละลายน้ำแข็ง}
\]

\[
Q_2=mL_f
\]

\[
\text{ช่วงที่ 3: ทำให้น้ำที่ละลายแล้วอุ่นขึ้น}
\]

\[
Q_3=m c_w(T_f-0)
\]

ดังนั้น

\[
Q_{\text{total}}=Q_1+Q_2+Q_3
\]

\subsection{อย่าเดาอุณหภูมิสุดท้ายก่อนตรวจสอบพลังงาน}

ในโจทย์น้ำแข็งกับน้ำร้อน อุณหภูมิสุดท้ายอาจเป็น

\[
T_f<0\,\si{\degreeCelsius}
\]

หรือ

\[
T_f=0\,\si{\degreeCelsius}
\]

หรือ

\[
T_f>0\,\si{\degreeCelsius}
\]

ต้องตรวจสอบว่าพลังงานที่มีเพียงพอสำหรับการหลอมละลายน้ำแข็งทั้งหมดหรือไม่

\begin{examplebox}{ตัวอย่าง: น้ำแข็งผสมกับน้ำ}
น้ำแข็งมวล $0.10\,\si{kg}$ ที่ $0\,\si{\degreeCelsius}$ ใส่ลงในน้ำมวล $0.40\,\si{kg}$ ที่ $30\,\si{\degreeCelsius}$ ในภาชนะฉนวน กำหนดให้

\[
c_w=4200\,\si{J/(kg.K)}
\]

\[
L_f=3.36\times10^5\,\si{J/kg}
\]

จงหาอุณหภูมิสุดท้าย

ถ้าน้ำเย็นลงถึง $0\,\si{\degreeCelsius}$ จะคายความร้อนได้

\[
Q_{\text{water}}=(0.40)(4200)(30)
\]

\[
Q_{\text{water}}=50400\,\si{J}
\]

พลังงานที่ใช้หลอมน้ำแข็งทั้งหมดคือ

\[
Q_{\text{melt}}=(0.10)(3.36\times10^5)
\]

\[
Q_{\text{melt}}=33600\,\si{J}
\]

มีพลังงานเหลือ

\[
Q_{\text{remain}}=50400-33600=16800\,\si{J}
\]

หลังน้ำแข็งละลายหมด มวลน้ำรวมคือ

\[
0.10+0.40=0.50\,\si{kg}
\]

พลังงานที่เหลือใช้เพิ่มอุณหภูมิน้ำทั้งหมด

\[
16800=(0.50)(4200)T_f
\]

ดังนั้น

\[
\boxed{T_f=8.0\,\si{\degreeCelsius}}
\]
\end{examplebox}

\begin{exercisebox}{แบบฝึกหัด 2: แคลอริเมตรีและการเปลี่ยนสถานะ}
\begin{enumerate}
    \item น้ำมวล $m$ ที่ $90\,\si{\degreeCelsius}$ ผสมกับน้ำมวล $3m$ ที่ $10\,\si{\degreeCelsius}$ ในภาชนะฉนวน จงหาอุณหภูมิสุดท้าย
    \item น้ำแข็งมวล $m$ ที่ $0\,\si{\degreeCelsius}$ ต้องใช้พลังงานเท่าไรจึงจะกลายเป็นน้ำที่ $20\,\si{\degreeCelsius}$
    \item น้ำแข็งมวล $20\,\si{g}$ ที่ $-10\,\si{\degreeCelsius}$ ใส่ลงในน้ำมวล $200\,\si{g}$ ที่ $10\,\si{\degreeCelsius}$ จงวิเคราะห์ว่าน้ำแข็งละลายหมดหรือไม่ โดยใช้ $c_{\text{ice}}=0.50\,\si{cal/(g.\degreeCelsius)}$, $c_w=1.0\,\si{cal/(g.\degreeCelsius)}$, และ $L_f=80\,\si{cal/g}$
    \item ไอน้ำมวล $m_s$ ที่ $100\,\si{\degreeCelsius}$ ควบแน่นและทำให้น้ำมวล $M$ ที่ $20\,\si{\degreeCelsius}$ มีอุณหภูมิสุดท้าย $60\,\si{\degreeCelsius}$ จงตั้งสมการหา $m_s$ โดยใช้ $L_v$ และ $c_w$
\end{enumerate}
\end{exercisebox}

% =========================================================
\section{การถ่ายโอนพลังงานความร้อน}

\subsection{สามกลไกหลัก}

พลังงานความร้อนถ่ายโอนได้สามแบบ

\begin{enumerate}
    \item การนำความร้อน
    \item การพาความร้อน
    \item การแผ่รังสีความร้อน
\end{enumerate}

สำหรับข้อสอบค่าย 1 หัวข้อที่ต้องคำนวณอย่างเป็นระบบที่สุดในบทนี้คือการนำความร้อน ส่วนการแผ่รังสีควรรู้ในระดับเสริมเพื่อรับมือโจทย์ประยุกต์

% =========================================================
\section{การนำความร้อน}

\subsection{ฟลักซ์ของการไหลของความร้อน}

ให้พลังงานความร้อนไหลผ่านแท่งวัตถุในแนวแกน $x$ ฟลักซ์ความร้อน $J$ คือพลังงานที่ไหลผ่านต่อหนึ่งหน่วยพื้นที่ต่อหนึ่งหน่วยเวลา

\[
J=\frac{1}{A}\frac{dQ}{dt}
\]

กฎการนำความร้อนของฟูเรียร์เขียนได้เป็น

\[
J=-K\frac{dT}{dx}
\]

เครื่องหมายลบหมายความว่าพลังงานความร้อนไหลจากบริเวณที่มีอุณหภูมิสูงกว่าไปยังบริเวณที่มีอุณหภูมิต่ำกว่า

\begin{tcolorbox}[title=ฟลักซ์ความร้อน]
\[
\boxed{J=-K\frac{dT}{dx}}
\]

สำหรับความชันอุณหภูมิประมาณคงตัว

\[
\boxed{J\approx-K\frac{\Delta T}{\Delta x}}
\]
\end{tcolorbox}

\subsection{อัตราการไหลของพลังงานความร้อนผ่านแท่ง}

ให้แท่งมีความยาว $L$ พื้นที่หน้าตัด $A$ และค่าสภาพนำความร้อน $K$

ปลายด้านร้อนมีอุณหภูมิ

\[
T_h
\]

ปลายด้านเย็นมีอุณหภูมิ

\[
T_c
\]

โดย

\[
T_h>T_c
\]

ขนาดของฟลักซ์คือ

\[
J=K\frac{T_h-T_c}{L}
\]

อัตราการไหลของพลังงานความร้อนคือ

\[
\mathcal{R}=\frac{dQ}{dt}=JA
\]

ดังนั้น

\[
\mathcal{R}
=
\frac{KA(T_h-T_c)}{L}
\]

\begin{tcolorbox}[title=อัตราการนำความร้อน]
\[
\boxed{
\mathcal{R}
=
\frac{dQ}{dt}
=
\frac{KA(T_h-T_c)}{L}
}
\]
\end{tcolorbox}

\subsection{ความต้านทานความร้อน}

สมการการนำความร้อนมีรูปคล้ายกฎของโอห์ม

\[
\mathcal{R}
=
\frac{T_h-T_c}{L/(KA)}
\]

นิยามความต้านทานความร้อนเป็น

\[
R_{\text{th}}=\frac{L}{KA}
\]

ดังนั้น

\[
\mathcal{R}=\frac{\Delta T}{R_{\text{th}}}
\]

\begin{tcolorbox}[title=ความต้านทานความร้อน]
\[
\boxed{R_{\text{th}}=\frac{L}{KA}}
\]

\[
\boxed{\mathcal{R}=\frac{\Delta T}{R_{\text{th}}}}
\]
\end{tcolorbox}

\subsection{แท่งหลายท่อนต่ออนุกรม}

ถ้าแท่งหลายท่อนต่อกันเป็นอนุกรม และระบบอยู่ในสภาวะคงตัว อัตราการไหลของพลังงานความร้อนผ่านทุกท่อนต้องเท่ากัน

\[
\mathcal{R}_1=\mathcal{R}_2=\cdots
\]

ความต้านทานความร้อนรวมคือ

\[
R_{\text{th,total}}
=
R_{\text{th},1}
+
R_{\text{th},2}
+\cdots
\]

ดังนั้น

\[
\mathcal{R}
=
\frac{T_h-T_c}{R_{\text{th,total}}}
\]

\begin{examplebox}{ตัวอย่าง: หาอุณหภูมิที่รอยต่อ}
แท่งสองท่อนทำด้วยวัสดุต่างกัน ต่ออนุกรมกัน ท่อนแรกมี

\[
L_1=L,\qquad A_1=A,\qquad K_1=K
\]

ท่อนที่สองมี

\[
L_2=2L,\qquad A_2=A,\qquad K_2=2K
\]

ปลายซ้ายมีอุณหภูมิ $T_1$ และปลายขวามีอุณหภูมิ $T_2$ โดย $T_1>T_2$ จงหาอุณหภูมิที่รอยต่อ $T_j$

ความต้านทานความร้อนของท่อนแรกคือ

\[
R_1=\frac{L}{KA}
\]

ท่อนที่สองคือ

\[
R_2=\frac{2L}{(2K)A}
\]

\[
R_2=\frac{L}{KA}
\]

ดังนั้น

\[
R_1=R_2
\]

เมื่ออยู่ในสภาวะคงตัว อัตราการไหลเท่ากัน และความต้านทานเท่ากัน จึงมีผลต่างอุณหภูมิของแต่ละท่อนเท่ากัน

\[
T_1-T_j=T_j-T_2
\]

ดังนั้น

\[
\boxed{T_j=\frac{T_1+T_2}{2}}
\]
\end{examplebox}

\subsection{แท่งเนื้อเดียวและกราฟอุณหภูมิ}

สำหรับแท่งเนื้อเดียว พื้นที่หน้าตัดคงตัว และไม่มีพลังงานความร้อนรั่วออกด้านข้าง ในสภาวะคงตัว

\[
\frac{dT}{dx}=\text{ค่าคงตัว}
\]

ดังนั้นกราฟ $T$ กับ $x$ เป็นเส้นตรง

ถ้าปลายซ้ายที่ $x=0$ มีอุณหภูมิ $T_1$ และปลายขวาที่ $x=L$ มีอุณหภูมิ $T_2$

\[
\boxed{
T(x)=T_1-\frac{x}{L}(T_1-T_2)
}
\]

\begin{exercisebox}{แบบฝึกหัด 3: การนำความร้อน}
\begin{enumerate}
    \item แท่งยาว $L$ พื้นที่หน้าตัด $A$ มีค่าสภาพนำความร้อน $K$ ปลายสองด้านต่างกัน $\Delta T$ จงหาอัตราการไหลของพลังงานความร้อน
    \item ถ้าความยาวของแท่งเพิ่มเป็นสองเท่า แต่พื้นที่หน้าตัดและผลต่างอุณหภูมิคงเดิม อัตราการไหลของพลังงานความร้อนเปลี่ยนเป็นกี่เท่า
    \item แท่งเนื้อเดียวปลายซ้าย $100\,\si{\degreeCelsius}$ ปลายขวา $20\,\si{\degreeCelsius}$ จุดหนึ่งอยู่ห่างจากปลายซ้ายเป็นหนึ่งในสี่ของความยาวแท่ง จุดนี้มีอุณหภูมิเท่าไร
    \item แท่งสองท่อนต่ออนุกรมกัน มีความต้านทานความร้อน $R$ และ $3R$ ถ้าปลายร้อนและปลายเย็นต่างกัน $80\,\si{K}$ ผลต่างอุณหภูมิคร่อมท่อนที่มีความต้านทาน $3R$ เท่าไร
\end{enumerate}
\end{exercisebox}

% =========================================================
\section{อ่านเสริมสำหรับโจทย์ประยุกต์: การแผ่รังสีความร้อน}

\subsection{กฎสเตฟาน--โบลทซ์มันน์}

วัตถุที่มีอุณหภูมิสัมบูรณ์ $T$ แผ่พลังงานความร้อนออกด้วยอัตรา

\[
P_{\text{rad}}=\varepsilon\sigma AT^4
\]

โดยที่

\[
\varepsilon=\text{สภาพเปล่งรังสี}
\]

\[
\sigma=\text{ค่าคงที่สเตฟาน--โบลทซ์มันน์}
\]

ถ้าต้องคิดพลังงานที่วัตถุรับกลับจากสิ่งแวดล้อมที่อุณหภูมิ $T_{\text{sur}}$

\[
P_{\text{net}}
=
\varepsilon\sigma A
\left(T^4-T_{\text{sur}}^4\right)
\]

\begin{tcolorbox}[title=หมายเหตุเรื่องขอบเขต]
หัวข้อหลักของบทนี้คือการนำความร้อนและการขยายตัวเชิงความร้อน แต่ข้อสอบบางปีใช้การนำความร้อนร่วมกับการแผ่รังสี จึงควรรู้สมการนี้ในฐานะเครื่องมือเสริม
\end{tcolorbox}

\subsection{สมดุลอัตราการไหลในสภาวะคงตัว}

ถ้าด้านหนึ่งของแผ่นรับพลังงานจากการนำความร้อน และอีกด้านคายพลังงานด้วยการแผ่รังสี เมื่อระบบเข้าสู่สภาวะคงตัว จะมี

\[
P_{\text{conduct}}=P_{\text{radiate}}
\]

เช่น

\[
\frac{KA(T_1-T_2)}{\ell}
=
\sigma AT_2^4
\]

ตัดพื้นที่ $A$

\[
\frac{K(T_1-T_2)}{\ell}
=
\sigma T_2^4
\]

% =========================================================
\section{วิธีเลือกหลักในการแก้โจทย์}

\begin{tcolorbox}[title=แผนการเลือกสมการ]
\begin{enumerate}
    \item ถ้าโจทย์ถามการเปลี่ยนขนาดของวัตถุเมื่ออุณหภูมิเปลี่ยน ใช้
    \[
    \Delta L=\alpha L_0\Delta T
    \]
    หรือ
    \[
    \Delta V=\beta V_0\Delta T
    \]

    \item ถ้าโจทย์ผสมวัตถุในภาชนะฉนวน ใช้
    \[
    \sum Q=0
    \]

    \item ถ้ามีการเปลี่ยนสถานะ ให้แบ่งช่วงและใช้
    \[
    Q=mL
    \]

    \item ถ้าโจทย์ถามพลังงานความร้อนไหลผ่านแท่งต่อหน่วยเวลา ใช้
    \[
    \mathcal{R}=\frac{KA\Delta T}{L}
    \]

    \item ถ้าแท่งหลายท่อนต่อกันในสภาวะคงตัว ใช้
    \[
    \mathcal{R}_1=\mathcal{R}_2
    \]
    หรือใช้ความต้านทานความร้อน
    \[
    R_{\text{th}}=\frac{L}{KA}
    \]
\end{enumerate}
\end{tcolorbox}

% =========================================================
\section{ชุดโจทย์รวมท้ายบท: ระดับ สอวน. ค่าย 1}

\subsection*{ระดับปรับพื้นฐาน}

\begin{enumerate}
    \item แผ่นโลหะมีรูวงกลมรัศมี $R$ และมีสัมประสิทธิ์การขยายตัวเชิงเส้น $\alpha$ เมื่ออุณหภูมิเพิ่มขึ้น $\Delta T$ จงหารัศมีใหม่ของรู โดยละพจน์อันดับสองขึ้นไป

    \item โลหะก้อนหนึ่งมีความจุความร้อนจำเพาะ $c_m$ มวล $m$ และอุณหภูมิเริ่มต้น $T_m$ ใส่ลงในน้ำมวล $M$ ที่อุณหภูมิ $T_w$ ในภาชนะฉนวน จงหาอุณหภูมิสมดุล $T_f$ ในกรณีที่ไม่มีการเปลี่ยนสถานะ

    \item แท่งนำความร้อนยาว $L$ พื้นที่หน้าตัด $A$ มีค่าสภาพนำความร้อน $K$ ปลายทั้งสองมีอุณหภูมิ $T_1>T_2$ จงหาอัตราการไหลของพลังงานความร้อนและฟลักซ์ความร้อน
\end{enumerate}

\subsection*{ระดับกลาง}

\begin{enumerate}
    \setcounter{enumi}{3}
    \item วงแหวนโลหะมีรัศมีภายใน $R$ ที่อุณหภูมิ $T_0$ ต้องการสวมเข้ากับแกนทรงกระบอกรัศมี $R+\delta$ โดย $\delta\ll R$ โลหะมีสัมประสิทธิ์การขยายตัวเชิงเส้น $\alpha$ จงหาอุณหภูมิที่ต้องเพิ่มขึ้นอย่างน้อย

    \item น้ำมวล $200\,\si{g}$ ที่ $40\,\si{\degreeCelsius}$ ผสมกับน้ำแข็งมวล $50\,\si{g}$ ที่ $-10\,\si{\degreeCelsius}$ ในภาชนะฉนวน กำหนดให้
    \[
    c_w=1.0\,\si{cal/(g.\degreeCelsius)}
    \]
    \[
    c_{\text{ice}}=0.50\,\si{cal/(g.\degreeCelsius)}
    \]
    \[
    L_f=80\,\si{cal/g}
    \]
    จงหาอุณหภูมิสุดท้าย

    \item แท่งสองท่อนต่ออนุกรมกัน ท่อนแรกมี $(L,A,K)$ และท่อนที่สองมี $(2L,3A,2K)$ ปลายซ้ายอุณหภูมิ $T_h$ และปลายขวาอุณหภูมิ $T_c$ จงหาอุณหภูมิที่รอยต่อ

    \item ลูกตุ้มอย่างง่ายมีความยาว $\ell$ และคาบ $T$ เส้นลวดมีสัมประสิทธิ์การขยายตัวเชิงเส้น $\alpha$ ถ้าอุณหภูมิเพิ่มขึ้น $\Delta T$ จงหาอัตราส่วน $T'/T$ แบบ exact และแบบประมาณอันดับหนึ่ง
\end{enumerate}

\subsection*{ระดับยาก}

\begin{enumerate}
    \setcounter{enumi}{7}
    \item แท่งวัสดุสองชนิดยาวเท่ากัน $L$ ถูกยึดติดกันตลอดแนว เมื่ออุณหภูมิเพิ่มขึ้น $\Delta T$ ระบบโค้งเป็นส่วนหนึ่งของวงกลม โดยรัศมีของชั้นนอกและชั้นในต่างกัน $d$ จงหาผลต่างของสัมประสิทธิ์การขยายตัวเชิงเส้นของวัสดุทั้งสอง

    \item ผนังประกอบด้วยชั้นวัสดุสองชั้นต่ออนุกรม ชั้นแรกมีความหนา $L$ ค่าสภาพนำความร้อน $K$ ชั้นที่สองมีความหนา $2L$ ค่าสภาพนำความร้อน $3K$ และพื้นที่หน้าตัดเท่ากัน ด้านในมีอุณหภูมิ $T_h$ ด้านนอกมีอุณหภูมิ $T_c$ จงหาอัตราการไหลของพลังงานความร้อนต่อหนึ่งหน่วยพื้นที่ และอุณหภูมิที่รอยต่อ

    \item น้ำมวล $M$ ที่อุณหภูมิ $T_0<100\,\si{\degreeCelsius}$ รับไอน้ำมวล $m$ ที่ $100\,\si{\degreeCelsius}$ จนกระทั่งอุณหภูมิสุดท้ายเป็น $T_f<100\,\si{\degreeCelsius}$ โดยภาชนะเป็นฉนวน จงหา $m/M$ ในรูปของ $c_w$, $L_v$, $T_0$, และ $T_f$

    \item แผ่นกันความร้อนหนา $\ell$ มีค่าสภาพนำความร้อน $K$ ด้านในมีอุณหภูมิ $T_1$ ด้านนอกมีอุณหภูมิ $T_2$ และแผ่รังสีออกสู่อวกาศโดยประมาณ จงแสดงว่า ถ้าให้
    \[
    x=\frac{T_2}{T_1}
    \]
    จะได้สมการ
    \[
    \left(\frac{\sigma\ell T_1^3}{K}\right)x^4+x=1
    \]

    \item กระเปาะบรรจุของเหลวปริมาตร $V$ ต่อกับท่อผอมพื้นที่หน้าตัด $a$ เมื่ออุณหภูมิเพิ่มขึ้น $\Delta T$ ของเหลวมีสัมประสิทธิ์การขยายตัวเชิงปริมาตร $\beta$ โดยไม่ต้องคิดการขยายตัวของภาชนะ จงหาระยะที่ระดับของเหลวเลื่อนขึ้นในท่อ
\end{enumerate}

% =========================================================
\section{เฉลยย่อท้ายบท}

\begin{enumerate}
    \item
    \[
    R'=R(1+\alpha\Delta T)
    \]

    \item ใช้
    \[
    mc_m(T_m-T_f)=Mc_w(T_f-T_w)
    \]
    ดังนั้น
    \[
    \boxed{
    T_f=
    \frac{mc_mT_m+Mc_wT_w}
    {mc_m+Mc_w}
    }
    \]

    \item
    \[
    \boxed{
    \mathcal{R}=\frac{KA(T_1-T_2)}{L}
    }
    \]
    \[
    \boxed{
    J=\frac{\mathcal{R}}{A}
    =
    \frac{K(T_1-T_2)}{L}
    }
    \]

    \item ต้องมี
    \[
    R(1+\alpha\Delta T)\geq R+\delta
    \]
    ดังนั้น
    \[
    \boxed{
    \Delta T_{\min}=\frac{\delta}{\alpha R}
    }
    \]

    \item น้ำคายพลังงานได้เมื่อเย็นลงถึง $0\,\si{\degreeCelsius}$
    \[
    Q_w=(200)(1.0)(40)=8000\,\si{cal}
    \]
    น้ำแข็งรับพลังงานเพื่ออุ่นถึง $0\,\si{\degreeCelsius}$
    \[
    Q_1=(50)(0.50)(10)=250\,\si{cal}
    \]
    ใช้หลอมละลายทั้งหมด
    \[
    Q_2=(50)(80)=4000\,\si{cal}
    \]
    พลังงานเหลือ
    \[
    Q_{\text{remain}}=8000-250-4000=3750\,\si{cal}
    \]
    มวลน้ำรวมคือ
    \[
    200+50=250\,\si{g}
    \]
    ดังนั้น
    \[
    3750=(250)(1.0)T_f
    \]
    \[
    \boxed{T_f=15\,\si{\degreeCelsius}}
    \]

    \item ความต้านทานความร้อนของท่อนแรกคือ
    \[
    R_1=\frac{L}{KA}
    \]
    ท่อนที่สองคือ
    \[
    R_2=\frac{2L}{(2K)(3A)}
    =
    \frac{L}{3KA}
    \]
    ดังนั้น
    \[
    \frac{T_h-T_j}{T_j-T_c}
    =
    \frac{R_1}{R_2}
    =
    3
    \]
    \[
    T_h-T_j=3T_j-3T_c
    \]
    \[
    \boxed{
    T_j=\frac{T_h+3T_c}{4}
    }
    \]

    \item
    \[
    \ell'=\ell(1+\alpha\Delta T)
    \]
    และ
    \[
    T=2\pi\sqrt{\frac{\ell}{g}}
    \]
    ดังนั้น
    \[
    \boxed{
    \frac{T'}{T}=\sqrt{1+\alpha\Delta T}
    }
    \]
    สำหรับ
    \[
    \alpha\Delta T\ll 1
    \]
    \[
    \boxed{
    \frac{T'}{T}
    \approx
    1+\frac{1}{2}\alpha\Delta T
    }
    \]

    \item ความยาวส่วนโค้งที่ต่างกันคือ
    \[
    \Delta L_{\text{diff}}=2\pi d
    \]
    อีกด้านหนึ่ง
    \[
    \Delta L_{\text{diff}}
    =
    L|\alpha_1-\alpha_2|\Delta T
    \]
    ดังนั้น
    \[
    \boxed{
    |\alpha_1-\alpha_2|
    =
    \frac{2\pi d}{L\Delta T}
    }
    \]

    \item ความต้านทานความร้อนต่อหนึ่งหน่วยพื้นที่คือ
    \[
    \frac{L}{K}
    +
    \frac{2L}{3K}
    =
    \frac{5L}{3K}
    \]
    ดังนั้นฟลักซ์คือ
    \[
    \boxed{
    J=
    \frac{T_h-T_c}{5L/(3K)}
    =
    \frac{3K(T_h-T_c)}{5L}
    }
    \]
    สำหรับรอยต่อ
    \[
    T_h-T_j
    =
    J\frac{L}{K}
    \]
    \[
    T_h-T_j
    =
    \frac{3}{5}(T_h-T_c)
    \]
    ดังนั้น
    \[
    \boxed{
    T_j=
    \frac{2T_h+3T_c}{5}
    }
    \]

    \item ไอน้ำคายพลังงานจากการควบแน่นและการเย็นลง
    \[
    Q_{\text{steam}}
    =
    mL_v
    +
    mc_w(100-T_f)
    \]
    น้ำเดิมรับพลังงาน
    \[
    Q_{\text{water}}
    =
    Mc_w(T_f-T_0)
    \]
    ตั้งให้เท่ากัน
    \[
    m\left[L_v+c_w(100-T_f)\right]
    =
    Mc_w(T_f-T_0)
    \]
    ดังนั้น
    \[
    \boxed{
    \frac{m}{M}
    =
    \frac{c_w(T_f-T_0)}
    {L_v+c_w(100-T_f)}
    }
    \]

    \item สภาวะคงตัวให้
    \[
    \frac{K(T_1-T_2)}{\ell}
    =
    \sigma T_2^4
    \]
    แทน
    \[
    T_2=xT_1
    \]
    จะได้
    \[
    \frac{KT_1(1-x)}{\ell}
    =
    \sigma x^4T_1^4
    \]
    จัดรูปเป็น
    \[
    \boxed{
    \left(\frac{\sigma\ell T_1^3}{K}\right)x^4+x=1
    }
    \]

    \item ปริมาตรเพิ่มคือ
    \[
    \Delta V=\beta V\Delta T
    \]
    ปริมาตรนี้เข้าไปอยู่ในท่อ
    \[
    a\Delta h=\Delta V
    \]
    ดังนั้น
    \[
    \boxed{
    \Delta h=\frac{\beta V\Delta T}{a}
    }
    \]
\end{enumerate}

% =========================================================
\section{สรุปสูตรสำคัญ}

\begin{tcolorbox}[title=สรุปบท: อุณหภูมิ ความร้อน และการขยายตัวเชิงความร้อน]
อุณหภูมิ

\[
T_{\si{K}}=T_{\si{\degreeCelsius}}+273.15
\]

การขยายตัวเชิงเส้น

\[
\Delta L=\alpha L_0\Delta T
\]

\[
L=L_0(1+\alpha\Delta T)
\]

การขยายตัวเชิงพื้นที่

\[
\Delta A\approx 2\alpha A_0\Delta T
\]

การขยายตัวเชิงปริมาตร

\[
\Delta V=\beta V_0\Delta T
\]

\[
\beta\approx 3\alpha
\]
\end{tcolorbox}

\begin{tcolorbox}[title=สรุปบท: อุณหภูมิ ความร้อน และการขยายตัวเชิงความร้อน ต่อ]
ความร้อนที่ทำให้อุณหภูมิเปลี่ยน

\[
Q=mc\Delta T
\]

ความร้อนแฝง

\[
Q=mL
\]

แคลอริเมตรีในภาชนะฉนวน

\[
\sum_i Q_i=0
\]

ฟลักซ์ความร้อน

\[
J=-K\frac{dT}{dx}
\]

อัตราการนำความร้อน

\[
\mathcal{R}=\frac{KA(T_h-T_c)}{L}
\]

ความต้านทานความร้อน

\[
R_{\text{th}}=\frac{L}{KA}
\]

การแผ่รังสีความร้อนสำหรับอ่านเสริม

\[
P_{\text{rad}}=\varepsilon\sigma AT^4
\]
\end{tcolorbox}

% =========================================================
\section{ข้อสอบจริง สอวน. ที่ใช้ความรู้บทความร้อนและการขยายตัวเชิงความร้อน}

\begin{tcolorbox}[title=คำแนะนำการใส่ภาพข้อสอบ]
ให้ตัดภาพข้อสอบจริงแล้วบันทึกตามชื่อไฟล์ที่กำหนดไว้ในแต่ละข้อ ถ้ายังไม่มีไฟล์ ภาพจะไม่แสดง แต่เอกสารยังคอมไพล์ได้

ข้อสอบปี 2567 ข้อ 15 ใช้การแผ่รังสีร่วมกับการนำความร้อน จึงจัดเป็นข้ออ่านเสริมระดับยาก
\end{tcolorbox}

% =========================================================
\subsection{ปี 2560 ตอนที่ 1 ข้อ 18: น้ำร้อนผสมกับน้ำมันในภาชนะฉนวน}

\vspace{1em}

\IfFileExists{img/posn60p1q18.png}{
\begin{center}
    \includegraphics[width=1\linewidth]{img/posn60p1q18.png}
\end{center}
}{
\begin{tcolorbox}[title=ตำแหน่งภาพที่แนะนำ]
ให้ตัดภาพข้อสอบปี 2560 ตอนที่ 1 ข้อ 18 แล้วบันทึกเป็น

\[
\texttt{img/posn60p1q18.png}
\]
\end{tcolorbox}
}

\begin{examplebox}{เฉลยละเอียด}
ให้ $T$ เป็นอุณหภูมิเริ่มต้นของน้ำมัน

น้ำมันมีมวล

\[
m_o=0.075\,\si{kg}
\]

และความจุความร้อนจำเพาะ

\[
c_o=1.4\times10^3\,\si{J/(kg.K)}
\]

น้ำมีมวล

\[
m_w=0.250\,\si{kg}
\]

อุณหภูมิเริ่มต้น

\[
80\,\si{\degreeCelsius}
\]

และ

\[
c_w=4.2\times10^3\,\si{J/(kg.K)}
\]

โจทย์บอกว่าอุณหภูมิสุดท้ายของน้ำมันเป็นสามเท่าของอุณหภูมิเดิม ดังนั้น

\[
T_f=3T
\]

ภาชนะเป็นฉนวน จึงใช้

\[
Q_{\text{water lost}}
=
Q_{\text{oil gained}}
\]

\[
m_wc_w(80-3T)
=
m_oc_o(3T-T)
\]

แทนค่า

\[
(0.250)(4.2\times10^3)(80-3T)
=
(0.075)(1.4\times10^3)(2T)
\]

\[
1050(80-3T)=210T
\]

\[
84000-3150T=210T
\]

\[
84000=3360T
\]

ดังนั้น

\[
\boxed{T=25\,\si{\degreeCelsius}}
\]
\end{examplebox}

% =========================================================
\subsection{ปี 2560 ตอนที่ 2 ข้อ 9: ทรงกระบอกและทรงกลมทองแดงขยายตัว}

\vspace{1em}

\IfFileExists{img/posn60p2q9.png}{
\begin{center}
    \includegraphics[width=1\linewidth]{img/posn60p2q9.png}
\end{center}
}{
\begin{tcolorbox}[title=ตำแหน่งภาพที่แนะนำ]
ให้ตัดภาพข้อสอบปี 2560 ตอนที่ 2 ข้อ 9 แล้วบันทึกเป็น

\[
\texttt{img/posn60p2q9.png}
\]
\end{tcolorbox}
}

\begin{examplebox}{เฉลยละเอียด}
ให้เส้นผ่านศูนย์กลางร่วมของทรงกระบอกและทรงกลมเป็น

\[
d
\]

และให้ความสูงของทรงกระบอกเป็น

\[
h
\]

ปริมาตรของทรงกระบอกคือ

\[
V_{\text{cyl}}
=
\pi\left(\frac{d}{2}\right)^2h
\]

ปริมาตรของทรงกลมคือ

\[
V_{\text{sphere}}
=
\frac{4}{3}\pi\left(\frac{d}{2}\right)^3
\]

โจทย์ให้ปริมาตรเริ่มต้นเท่ากัน

\[
\pi\left(\frac{d}{2}\right)^2h
=
\frac{4}{3}\pi\left(\frac{d}{2}\right)^3
\]

ตัดพจน์ร่วม

\[
h=\frac{2}{3}d
\]

เพราะวัตถุทั้งสองทำจากทองแดงเหมือนกันและได้รับความร้อนเท่ากัน โดยมวลเท่ากันจากปริมาตรที่เท่ากัน จึงมีผลต่างอุณหภูมิเท่ากัน และทุกระยะเชิงเส้นขยายด้วยสัดส่วนเดียวกัน

\[
\frac{\Delta h}{h}
=
\frac{\Delta d}{d}
\]

ดังนั้น

\[
\frac{\Delta h}{\Delta d}
=
\frac{h}{d}
\]

\[
\boxed{
\frac{\Delta h}{\Delta d}
=
\frac{2}{3}
}
\]
\end{examplebox}

% =========================================================
\subsection{ปี 2561 ข้อ 18: แถบโลหะสองชนิดโค้งเป็นวงกลม}

\vspace{1em}

\IfFileExists{img/posn61q18.png}{
\begin{center}
    \includegraphics[width=1\linewidth]{img/posn61q18.png}
\end{center}
}{
\begin{tcolorbox}[title=ตำแหน่งภาพที่แนะนำ]
ให้ตัดภาพข้อสอบปี 2561 ข้อ 18 แล้วบันทึกเป็น

\[
\texttt{img/posn61q18.png}
\]
\end{tcolorbox}
}

\begin{examplebox}{เฉลยละเอียด}
แท่งวัสดุสองชนิดยึดติดกัน เดิมมีความยาวเท่ากัน

\[
L=20\pi\,\si{cm}
\]

เมื่ออุณหภูมิเพิ่มขึ้น

\[
\Delta T=80\,\si{K}
\]

ระบบโค้งเป็นวงกลม โดยรัศมีของวัสดุสองชั้นต่างกัน

\[
\Delta r=0.20\,\si{mm}
\]

เปลี่ยนหน่วย

\[
\Delta r=0.020\,\si{cm}
\]

ผลต่างของเส้นรอบวงคือ

\[
\Delta L_{\text{diff}}=2\pi\Delta r
\]

อีกด้านหนึ่ง ผลต่างของการขยายความยาวคือ

\[
\Delta L_{\text{diff}}
=
L|\alpha_1-\alpha_2|\Delta T
\]

ดังนั้น

\[
|\alpha_1-\alpha_2|
=
\frac{2\pi\Delta r}{L\Delta T}
\]

แทนค่า

\[
|\alpha_1-\alpha_2|
=
\frac{2\pi(0.020)}
{(20\pi)(80)}
\]

\[
\boxed{
|\alpha_1-\alpha_2|
=
2.5\times10^{-5}\,\si{K^{-1}}
}
\]
\end{examplebox}

% =========================================================
\subsection{ปี 2561 ตอนที่ 2 ข้อ 9: โลหะร้อน น้ำ และน้ำแข็ง}

\vspace{1em}

\IfFileExists{img/posn61p2q9.png}{
\begin{center}
    \includegraphics[width=1\linewidth]{img/posn61p2q9.png}
\end{center}
}{
\begin{tcolorbox}[title=ตำแหน่งภาพที่แนะนำ]
ให้ตัดภาพข้อสอบปี 2561 ตอนที่ 2 ข้อ 9 แล้วบันทึกเป็น

\[
\texttt{img/posn61p2q9.png}
\]
\end{tcolorbox}
}

\begin{examplebox}{เฉลยละเอียด}
ภาชนะฉนวนมีน้ำ

\[
m_w=300\,\si{g}
\]

และน้ำแข็ง

\[
m_i=200\,\si{g}
\]

ที่

\[
0\,\si{\degreeCelsius}
\]

ใส่โลหะมวล

\[
m_m=1000\,\si{g}
\]

ลงไป แล้วพบว่าอุณหภูมิสุดท้ายเป็น

\[
T_f=10\,\si{\degreeCelsius}
\]

กำหนดให้

\[
c_w=4.2\,\si{J/(g.K)}
\]

\[
L_f=335\,\si{J/g}
\]

\[
c_m=2.0\,\si{J/(g.K)}
\]

พลังงานที่ใช้หลอมน้ำแข็งคือ

\[
Q_1=m_iL_f
\]

\[
Q_1=(200)(335)
\]

\[
Q_1=67000\,\si{J}
\]

หลังน้ำแข็งละลาย น้ำทั้งหมดมีมวล

\[
300+200=500\,\si{g}
\]

พลังงานที่ใช้เพิ่มอุณหภูมิน้ำทั้งหมดจาก $0$ ถึง $10\,\si{\degreeCelsius}$ คือ

\[
Q_2=(500)(4.2)(10)
\]

\[
Q_2=21000\,\si{J}
\]

พลังงานรวมที่น้ำและน้ำแข็งรับคือ

\[
Q_{\text{gain}}=67000+21000
\]

\[
Q_{\text{gain}}=88000\,\si{J}
\]

ให้ $T_m$ เป็นอุณหภูมิเริ่มต้นของโลหะ โลหะคายพลังงาน

\[
Q_{\text{metal}}
=
(1000)(2.0)(T_m-10)
\]

ใช้อนุรักษ์พลังงาน

\[
(1000)(2.0)(T_m-10)=88000
\]

\[
T_m-10=44
\]

ดังนั้น

\[
\boxed{
T_m=54\,\si{\degreeCelsius}
}
\]
\end{examplebox}

% =========================================================
\subsection{ปี 2562 ข้อ 7: การกระจายอุณหภูมิในแท่งเนื้อเดียว}

\vspace{1em}

\IfFileExists{img/posn62q7.png}{
\begin{center}
    \includegraphics[width=1\linewidth]{img/posn62q7.png}
\end{center}
}{
\begin{tcolorbox}[title=ตำแหน่งภาพที่แนะนำ]
ให้ตัดภาพข้อสอบปี 2562 ข้อ 7 แล้วบันทึกเป็น

\[
\texttt{img/posn62q7.png}
\]
\end{tcolorbox}
}

\begin{examplebox}{เฉลยละเอียด}
แท่งโลหะเนื้อเดียวมีความยาว

\[
L
\]

พื้นที่หน้าตัดคงตัว และด้านข้างหุ้มฉนวน ปลาย $A$ มีอุณหภูมิ

\[
T_1
\]

ปลาย $B$ มีอุณหภูมิ

\[
T_2
\]

โดย

\[
T_1>T_2
\]

ในสภาวะคงตัว สำหรับแท่งเนื้อเดียว ความชันอุณหภูมิคงตัว ดังนั้นกราฟ $T$ กับ $x$ เป็นเส้นตรง

ที่

\[
x=0
\]

มี

\[
T=T_1
\]

และที่

\[
x=L
\]

มี

\[
T=T_2
\]

ดังนั้นที่ระยะ $x$ จากปลาย $A$

\[
\boxed{
T(x)
=
T_1-\frac{x}{L}(T_1-T_2)
}
\]
\end{examplebox}

% =========================================================
\subsection{ปี 2562 ข้อ 8: คาบลูกตุ้มเมื่อเส้นลวดขยายตัว}

\vspace{1em}

\IfFileExists{img/posn62q8.png}{
\begin{center}
    \includegraphics[width=1\linewidth]{img/posn62q8.png}
\end{center}
}{
\begin{tcolorbox}[title=ตำแหน่งภาพที่แนะนำ]
ให้ตัดภาพข้อสอบปี 2562 ข้อ 8 แล้วบันทึกเป็น

\[
\texttt{img/posn62q8.png}
\]
\end{tcolorbox}
}

\begin{examplebox}{เฉลยละเอียด}
คาบของลูกตุ้มอย่างง่ายคือ

\[
T=2\pi\sqrt{\frac{\ell}{g}}
\]

เมื่ออุณหภูมิเพิ่มขึ้น $\Delta T$ ความยาวลวดใหม่คือ

\[
\ell'=\ell(1+\alpha\Delta T)
\]

ดังนั้น

\[
\frac{T'}{T}
=
\sqrt{\frac{\ell'}{\ell}}
\]

\[
\boxed{
\frac{T'}{T}
=
\sqrt{1+\alpha\Delta T}
}
\]

ถ้า

\[
\alpha\Delta T\ll1
\]

ใช้

\[
\sqrt{1+\epsilon}\approx1+\frac{1}{2}\epsilon
\]

จึงได้

\[
\boxed{
\frac{T'}{T}
\approx
1+\frac{1}{2}\alpha\Delta T
}
\]
\end{examplebox}

% =========================================================
\subsection{อ่านเสริม: ข้อสอบศูนย์ สอวน. ปี 2563 ข้อ 8: อุณหภูมิในท่อนทองแดง}

\vspace{1em}

\IfFileExists{img/posn63centerq8.png}{
\begin{center}
    \includegraphics[width=1\linewidth]{img/posn63centerq8.png}
\end{center}
}{
\begin{tcolorbox}[title=ตำแหน่งภาพที่แนะนำ]
ให้ตัดภาพข้อสอบศูนย์ สอวน. ปี 2563 ข้อ 8 แล้วบันทึกเป็น

\[
\texttt{img/posn63centerq8.png}
\]
\end{tcolorbox}
}

\begin{examplebox}{เฉลยละเอียด}
ท่อนทองแดง $AB$ หุ้มฉนวนด้านข้าง ปลาย $A$ มีอุณหภูมิ

\[
50\,\si{\degreeCelsius}
\]

ปลาย $B$ มีอุณหภูมิ

\[
-10\,\si{\degreeCelsius}
\]

จุด $P$ อยู่ห่างจากปลาย $A$ เป็น

\[
\frac{1}{4}AB
\]

ในสภาวะคงตัว อุณหภูมิเปลี่ยนเป็นเส้นตรงตามตำแหน่ง

\[
T(x)=T_A-\frac{x}{L}(T_A-T_B)
\]

แทน

\[
\frac{x}{L}=\frac{1}{4}
\]

\[
T_P
=
50-\frac{1}{4}\left[50-(-10)\right]
\]

\[
T_P=50-15
\]

ดังนั้น

\[
\boxed{
T_P=35\,\si{\degreeCelsius}
}
\]
\end{examplebox}

% =========================================================
\subsection{อ่านเสริม: ข้อสอบศูนย์ สอวน. ปี 2563 ข้อ 9: น้ำแข็งผสมกับน้ำ}

\vspace{1em}

\IfFileExists{img/posn63centerq9.png}{
\begin{center}
    \includegraphics[width=1\linewidth]{img/posn63centerq9.png}
\end{center}
}{
\begin{tcolorbox}[title=ตำแหน่งภาพที่แนะนำ]
ให้ตัดภาพข้อสอบศูนย์ สอวน. ปี 2563 ข้อ 9 แล้วบันทึกเป็น

\[
\texttt{img/posn63centerq9.png}
\]
\end{tcolorbox}
}

\begin{examplebox}{เฉลยละเอียด}
น้ำแข็งมวล

\[
10\,\si{g}
\]

เริ่มที่

\[
-10\,\si{\degreeCelsius}
\]

ใส่ลงในน้ำมวล

\[
200\,\si{g}
\]

ที่

\[
10\,\si{\degreeCelsius}
\]

กำหนดให้

\[
c_{\text{ice}}=0.50\,\si{cal/(g.\degreeCelsius)}
\]

\[
c_w=1.0\,\si{cal/(g.\degreeCelsius)}
\]

\[
L_f=80.0\,\si{cal/g}
\]

น้ำคายพลังงานได้เมื่อเย็นลงถึง $0\,\si{\degreeCelsius}$

\[
Q_w=(200)(1.0)(10)
\]

\[
Q_w=2000\,\si{cal}
\]

พลังงานที่ใช้เพิ่มอุณหภูมิน้ำแข็งถึง $0\,\si{\degreeCelsius}$ คือ

\[
Q_1=(10)(0.50)(10)
\]

\[
Q_1=50\,\si{cal}
\]

พลังงานที่ใช้หลอมน้ำแข็งคือ

\[
Q_2=(10)(80.0)
\]

\[
Q_2=800\,\si{cal}
\]

พลังงานเหลือ

\[
Q_{\text{remain}}=2000-50-800
\]

\[
Q_{\text{remain}}=1150\,\si{cal}
\]

น้ำทั้งหมดหลังละลายมีมวล

\[
200+10=210\,\si{g}
\]

ดังนั้น

\[
1150=(210)(1.0)T_f
\]

\[
\boxed{
T_f\approx5.48\,\si{\degreeCelsius}
}
\]
\end{examplebox}

% =========================================================
\subsection{อ่านเสริม: ข้อสอบศูนย์ สอวน. ปี 2563 ข้อ 30: ระดับปรอทในท่อผอม}

\vspace{1em}

\IfFileExists{img/posn63centerq30.png}{
\begin{center}
    \includegraphics[width=1\linewidth]{img/posn63centerq30.png}
\end{center}
}{
\begin{tcolorbox}[title=ตำแหน่งภาพที่แนะนำ]
ให้ตัดภาพข้อสอบศูนย์ สอวน. ปี 2563 ข้อ 30 แล้วบันทึกเป็น

\[
\texttt{img/posn63centerq30.png}
\]
\end{tcolorbox}
}

\begin{examplebox}{เฉลยละเอียด}
กระเปาะใหญ่บรรจุปรอทปริมาตร

\[
V
\]

ต่อกับท่อผอมพื้นที่หน้าตัด

\[
a
\]

ให้สัมประสิทธิ์การขยายตัวเชิงปริมาตรของปรอทเป็น

\[
\gamma
\]

เมื่ออุณหภูมิเพิ่มขึ้น

\[
1\,\si{\degreeCelsius}
\]

ปริมาตรปรอทเพิ่มขึ้น

\[
\Delta V=\gamma V
\]

ปริมาตรที่เพิ่มเข้าไปอยู่ในท่อผอม

\[
\Delta V=a\Delta h
\]

ดังนั้น

\[
a\Delta h=\gamma V
\]

\[
\boxed{
\Delta h=\frac{\gamma V}{a}
}
\]
\end{examplebox}

% =========================================================
\subsection{อ่านเสริมระดับยาก: ปี 2567 ข้อ 15 การนำความร้อนร่วมกับการแผ่รังสี}

\vspace{1em}

\IfFileExists{img/posn67q15.png}{
\begin{center}
    \includegraphics[width=1\linewidth]{img/posn67q15.png}
\end{center}
}{
\begin{tcolorbox}[title=ตำแหน่งภาพที่แนะนำ]
ให้ตัดภาพข้อสอบปี 2567 ข้อ 15 แล้วบันทึกเป็น

\[
\texttt{img/posn67q15.png}
\]
\end{tcolorbox}
}

\begin{examplebox}{เฉลยละเอียด}
ให้แผ่นกันความร้อนมีความหนา

\[
\ell=0.10\,\si{m}
\]

ค่าสภาพนำความร้อน

\[
K=0.57\,\si{W/(m.K)}
\]

ผิวด้านในมีอุณหภูมิ

\[
T_1=1000\,\si{K}
\]

และผิวด้านนอกมีอุณหภูมิ

\[
T_2
\]

ผิวด้านนอกแผ่รังสีสู่อวกาศโดยประมาณ กำหนดให้

\[
\sigma=5.7\times10^{-8}\,\si{W/(m^2.K^4)}
\]

ในสภาวะคงตัว อัตราการนำพลังงานต่อพื้นที่เท่ากับอัตราการแผ่รังสีต่อพื้นที่

\[
\frac{K(T_1-T_2)}{\ell}
=
\sigma T_2^4
\]

ให้

\[
x=\frac{T_2}{T_1}
\]

ดังนั้น

\[
T_2=xT_1
\]

แทนในสมการ

\[
\frac{KT_1(1-x)}{\ell}
=
\sigma x^4T_1^4
\]

จัดรูป

\[
\left(
\frac{\sigma\ell T_1^3}{K}
\right)x^4+x=1
\]

แทนค่า

\[
\frac{\sigma\ell T_1^3}{K}
=
\frac{(5.7\times10^{-8})(0.10)(1000)^3}{0.57}
\]

\[
\frac{\sigma\ell T_1^3}{K}=10
\]

ดังนั้น

\[
10x^4+x=1
\]

โจทย์ให้รากเป็น

\[
x=0.478
\]

จึงได้

\[
T_2=xT_1
\]

\[
T_2=(0.478)(1000)
\]

\[
\boxed{
T_2=478\,\si{K}
}
\]
\end{examplebox}

% =========================================================
\subsection{ปี 2568 ข้อ 9: อัตราการไหลของพลังงานผ่านท่อนำความร้อน}

\vspace{1em}

\IfFileExists{img/posn68q9.png}{
\begin{center}
    \includegraphics[width=1\linewidth]{img/posn68q9.png}
\end{center}
}{
\begin{tcolorbox}[title=ตำแหน่งภาพที่แนะนำ]
ให้ตัดภาพข้อสอบปี 2568 ข้อ 9 แล้วบันทึกเป็น

\[
\texttt{img/posn68q9.png}
\]
\end{tcolorbox}
}

\begin{examplebox}{เฉลยละเอียด}
ท่อนำความร้อนมีความยาว

\[
L
\]

พื้นที่หน้าตัด

\[
A
\]

และค่าสภาพนำความร้อน

\[
K
\]

ปลายซ้ายอยู่ที่อุณหภูมิ

\[
T_1
\]

ปลายขวาอยู่ที่อุณหภูมิ

\[
T_2
\]

โดย

\[
T_1>T_2
\]

ขนาดของฟลักซ์ความร้อนคือ

\[
J=
K\frac{T_1-T_2}{L}
\]

อัตราที่พลังงานความร้อนไหลผ่านพื้นที่หน้าตัด $A$ คือ

\[
\mathcal{R}=JA
\]

ดังนั้น

\[
\boxed{
\mathcal{R}
=
\frac{KA(T_1-T_2)}{L}
}
\]
\end{examplebox}

% =========================================================
\subsection{ปี 2568 ข้อ 10: อุณหภูมิที่รอยต่อของท่อนำความร้อน}

\vspace{1em}

\IfFileExists{img/posn68q10.png}{
\begin{center}
    \includegraphics[width=1\linewidth]{img/posn68q10.png}
\end{center}
}{
\begin{tcolorbox}[title=ตำแหน่งภาพที่แนะนำ]
ให้ตัดภาพข้อสอบปี 2568 ข้อ 10 แล้วบันทึกเป็น

\[
\texttt{img/posn68q10.png}
\]
\end{tcolorbox}
}

\begin{examplebox}{เฉลยละเอียด}
ท่อนำความร้อนสองท่อนทำด้วยสารชนิดเดียวกัน

ท่อนแรกมี

\[
L_1=L
\]

\[
A_1=A
\]

ท่อนที่สองมี

\[
L_2=2L
\]

\[
A_2=2A
\]

ค่าสภาพนำความร้อนเท่ากันคือ

\[
K
\]

ความต้านทานความร้อนของท่อนแรกคือ

\[
R_1=\frac{L}{KA}
\]

ของท่อนที่สองคือ

\[
R_2=\frac{2L}{K(2A)}
\]

\[
R_2=\frac{L}{KA}
\]

ดังนั้น

\[
R_1=R_2
\]

เมื่ออยู่ในสภาวะคงตัว อัตราการไหลของพลังงานผ่านสองท่อนเท่ากัน และเมื่อความต้านทานเท่ากัน ผลต่างอุณหภูมิคร่อมแต่ละท่อนจึงเท่ากัน

ให้ $T_j$ เป็นอุณหภูมิที่รอยต่อ

\[
T_1-T_j=T_j-T_2
\]

ดังนั้น

\[
\boxed{
T_j=\frac{T_1+T_2}{2}
}
\]
\end{examplebox}